\section{Introduction}

Ce projet constitue un exemple prêt à remplacer. Le corps du texte est composé en Times New Roman 12 points, à interligne 1,5, sans alinéa et avec un alignement justifié. Une source peut être citée entre parenthèses \parencite[33--40]{gaudreau2011} ou de manière narrative avec la commande \textcite{loiseau2026}.

La structure annoncée dans l'introduction doit correspondre à celle du développement et de la table des matières. L'introduction présente habituellement le sujet amené, le sujet posé et le sujet divisé.

\section{Développement}

\subsection{Une première sous-section}

Une citation directe de trois lignes ou moins demeure dans le paragraphe, entre guillemets français : \enquote{insérez ici une courte citation correctement référencée} \parencite[18]{loiseau2026}.

Une citation directe de plus de trois lignes se place sans guillemets, à interligne simple et avec un retrait de 2,5 cm de chaque côté :

\begin{quote}
Insérez ici une citation de plus de trois lignes. Le texte est automatiquement placé à interligne simple, justifié, et en retrait de 2,5 centimètres à gauche comme à droite. Ajoutez ensuite la référence auteur-date et la page exacte \parencite[19]{loiseau2026}.
\end{quote}

Une note infrapaginale explicative est automatiquement composée dans la même police, en taille 10, à interligne simple et alignée à gauche.\footnote{Ceci est un exemple de note explicative; les sources demeurent citées dans le texte selon le système auteur-date APA.}

\subsection{Tableaux et figures}

Le Tableau~\ref{tab:exemple} illustre la numérotation liée à la section. Chaque tableau ou figure doit être mentionné dans le texte et accompagné d'une source.

\begin{table}[H]
  \centering
  \caption{Exemple de tableau}
  \label{tab:exemple}
  \begin{tabular}{lcc}
    \hline
    Catégorie & Valeur A & Valeur B \\
    \hline
    Exemple 1 & 10 & 20 \\
    Exemple 2 & 15 & 25 \\
    \hline
  \end{tabular}
\end{table}

{\fontsize{10}{12}\selectfont\raggedright Source : Création de l'autrice ou de l'auteur.\par}

\section{Conclusion}

La conclusion rappelle la problématique, la méthode et les principaux résultats. Elle présente ensuite les conclusions ou recommandations, propose une critique réflexive de la démarche et se termine par une ouverture pertinente.
